%% Question bank chapter file - generated by bank_tools.py build (do not edit while the bank is still being built:
%% it is regenerated from the per-test data in Question Bank\data). Every question is one bankq block:
%%   \begin{bankq}[id=..., source={...}, paper=.., qno=.., type=..., subtopic={...}, answer={...}, answer_src=key|solved,
%%                 working={...}, difficulty=1-10, diff_parts={K S I C F}, diff_note={...},
%%                 group=<passage id or empty>, also_in={duplicate copies in other tests}]
%% difficulty = min(10, K+S+I+C+F-1): K concepts 1-3, S steps 1-4, I insight 0-2, C computation 0-2, F format 0-1
%%   ...question LaTeX, options, figures (\bfig / \bfigi -> img/)...
%%   \end{bankq}
%% Passages shared by several questions are bankpassage blocks; their questions carry group=<passage id>.
\documentclass[10pt]{article}
\usepackage[a4paper,margin=1.8cm]{geometry}
\usepackage[T1]{fontenc}
\usepackage{amsmath,amssymb,mathtools}
\usepackage{graphicx}
\usepackage[export]{adjustbox}
\usepackage{xcolor,needspace,keyval}
\usepackage{enumitem}
\usepackage{booktabs,array,tabularx}
\usepackage[hidelinks]{hyperref}
\setlength{\parindent}{0pt}
\setlength{\parskip}{3pt}
\makeatletter
\define@key{bq}{id}{\def\bq@id{#1}}\define@key{bq}{source}{\def\bq@src{#1}}\define@key{bq}{paper}{\def\bq@paper{#1}}
\define@key{bq}{qno}{\def\bq@qno{#1}}\define@key{bq}{type}{\def\bq@type{#1}}\define@key{bq}{subtopic}{\def\bq@sub{#1}}
\define@key{bq}{answer}{\def\bq@ans{#1}}\define@key{bq}{answer_src}{\def\bq@asrc{#1}}\define@key{bq}{working}{\def\bq@work{#1}}
\define@key{bq}{group}{\def\bq@group{#1}}\define@key{bq}{also_in}{\def\bq@also{#1}}
\define@key{bq}{difficulty}{\def\bq@diff{#1}}\define@key{bq}{diff_parts}{\def\bq@dparts{#1}}\define@key{bq}{diff_note}{\def\bq@dnote{#1}}
\newenvironment{bankq}[1][]{\def\bq@work{}\def\bq@also{}\def\bq@group{}\def\bq@diff{?}\def\bq@dparts{}\def\bq@dnote{}\setkeys{bq}{#1}%
  \par\needspace{4\baselineskip}\medskip\noindent{\small\textbf{\bq@id}\hfill\color{gray}[\bq@src, Paper \bq@paper, Q\bq@qno]
  \bq@type\ $\cdot$ \bq@sub\ $\cdot$ \textbf{D\bq@diff}}\par\nopagebreak\relax}
  {\par{\small\color{teal}Answer: \textbf{\bq@ans}\ifx\bq@work\@empty\else\ (solved: \bq@work)\fi
  \ifx\bq@also\@empty\else\par\color{gray}Also in: \bq@also\fi}\par}
\define@key{bp}{id}{\def\bp@id{#1}}\define@key{bp}{source}{\def\bp@src{#1}}\define@key{bp}{questions}{\def\bp@qs{#1}}
\newenvironment{bankpassage}[1][]{\setkeys{bp}{#1}\par\needspace{4\baselineskip}\bigskip\noindent
  {\small\textbf{Passage \bp@id}\hfill\color{gray}[\bp@src; questions \bp@qs]}\par\nopagebreak\relax\itshape}{\par}
\makeatother
\newcommand{\bfig}[2]{\par\smallskip{\centering\includegraphics[width=#1pt,max width=0.62\linewidth,max height=0.32\textheight]{img/#2}\par}\smallskip}
\newcommand{\bfigi}[2]{\includegraphics[width=#1pt,max width=\linewidth,max height=0.32\textheight,valign=c]{img/#2}}
\newcommand{\bmissing}[1]{\textit{[figure missing in the source conversion: \detokenize{#1}]}}
\newenvironment{qoptions}{\begin{itemize}[leftmargin=1.8em,itemsep=1pt,topsep=2pt,label={}]}{\end{itemize}}
\newcommand{\opt}[2]{\item[(#1)] #2}
\providecommand{\fwstart}{}\providecommand{\fwend}{}\providecommand{\qfullsplit}{}
\newcolumntype{L}[1]{>{\raggedright\arraybackslash}p{#1}}
\newcolumntype{C}[1]{>{\centering\arraybackslash}p{#1}}
\newcolumntype{R}[1]{>{\raggedleft\arraybackslash}p{#1}}
\providecommand{\sectionhdr}[1]{}\providecommand{\subsectionhdr}[1]{}


\begin{document}
{\Large\bfseries Thermochemistry --- Question Bank}\par\smallskip
{\small Chemistry chapter C05. Questions from all converted test papers, grouped by source series. Answers: printed key where it matches the paper, otherwise solved (answer\_src=solved, with working). Exact repeats across papers appear once, with also\_in listing the other copies.}\par\medskip\hrule

\section*{Score-I papers}
\begin{bankq}[id=SI12-P1-CHEM-S2-Q03, source={Score-I - Test 12 (09 Mar 2022)}, paper=1, qno=3, type=numerical, subtopic={Enthalpies}, answer={-242}, answer_src=solved, working={$\Delta U=-(8.75\times0.312)\times90=-245.7$ kJ/mol; $\Delta n_g=2-\tfrac12=1.5$; $\Delta H=\Delta U+1.5RT\approx-242$ kJ/mol}, difficulty=3, diff_parts={K1 S2 I0 C0 F1}, diff_note={concepts: enthalpy of combustion from a bomb calorimeter | steps: heat evolved per mole gives delta U; delta ng = 1.5; delta H = delta U + delta ng RT}]
When 1.0 g of oxalic acid (H$_2$C$_2$O$_4$) is burned in a bomb calorimeter whose heat capacity 8.75 kJ/K, the temperature increases by 0.312 K. The enthalpy of combustion of oxalic acid 27$^\circ$C is (kJ/mole) ($R=8.3$ J/mole-K)
\end{bankq}

\section*{Score Advanced papers}
\begin{bankq}[id=SA01-P2-CHEM-S2-Q12, source={Score Advanced - Test 01 (22 Sep 2020)}, paper=2, qno=12, type=single, subtopic={Enthalpies}, answer={D}, answer_src=key, difficulty=7, diff_parts={K2 S2 I2 C2 F0}, diff_note={concepts: enthalpy from formation data, enthalpies | steps: bond relations; reaction; value}]
Given the following data\par
$\Delta_f\mathrm H^\circ(\mathrm{CH_4},g)=-80\,\mathrm{kJ\,mol^{-1}}$\par
$\Delta_f\mathrm H^\circ(\mathrm{C_2H_6},g)=-85\,\mathrm{kJ\,mol^{-1}}$\par
$\Delta_f\mathrm H^\circ(\mathrm{C_2H_4},g)=50\,\mathrm{kJ\,mol^{-1}}$\par
Enthalpy of sublimation of Carbon (graphite) = 720 kJ mol$^{-1}$\par
Enthalpy of atomisation of H = 215 kJ mol$^{-1}$\par
\noindent{\footnotesize\setlength{\tabcolsep}{2pt}\begin{tabular}{|L{.7cm}|L{3.4cm}|L{.7cm}|L{2.2cm}|}\hline
\multicolumn{2}{|c|}{\parbox{4.0cm}{\centering List -- I (Thermodynamic data)}} & \multicolumn{2}{c|}{\parbox{2.6cm}{\centering List -- II (Values with respect to magnitude)}}\\\hline
(P) & Bond enthalpy of C--H bond (per mole) & (1) & 325 kJ\\\hline
(Q) & Bond enthalpy of C--C bond (per mole) & (2) & 600 kJ\\\hline
(R) & Bond enthalpy of C=C bond (per mole) & (3) & 415 kJ\\\hline
(S) & $\Delta$H$^\circ$, for the reaction : $\mathrm{nC_2H_4(g)}\longrightarrow$ $\mathrm{(-CH-CH_2-)_n(g)}$ (for 10 moles of ethene) & (4) & 590 kJ\\\hline
\end{tabular}}\par
The correct option is -
\begin{qoptions}
\opt{A}{P $\to$ 3; Q $\to$ 4; R $\to$ 1; S $\to$ 2}
\opt{B}{P $\to$ 3; Q $\to$ 4; R $\to$ 2; S $\to$ 1}
\opt{C}{P $\to$ 1; Q $\to$ 3; R $\to$ 2; S $\to$ 4}
\opt{D}{P $\to$ 3; Q $\to$ 1; R $\to$ 4; S $\to$ 2}
\end{qoptions}
\end{bankq}
\begin{bankq}[id=SA03-P2-CHEM-S2-Q05, source={Score Advanced - Test 03 (27 Jun 2021)}, paper=2, qno=5, type=single, subtopic={Enthalpies}, answer={-38}, answer_src=solved, working={-50-45+57=-38 kJ}, difficulty=2, diff_parts={K1 S1 I1 C0 F0}, diff_note={concepts: enthalpy by Hess law | steps: combine the reactions}]
\par\noindent{\footnotesize\begin{tabular}{@{}l@{\ }l@{}}
HCl + BOH $\rightarrow$ BCl + H$_2$O & $\Delta$H = $-$57 kJ\\
CH$_3$COOH + BOH $\rightarrow$ CH$_3$COOB + H$_2$O & $\Delta$H = $-$50 kJ\\
HCl + DOH $\rightarrow$ DCl + H$_2$O & $\Delta$H = $-$45 kJ
\end{tabular}}\par
then $\Delta$H (in kJ) for CH$_3$COOH + DOH $\rightarrow$ CH$_3$COOD + H$_2$O, is equal to :-
\end{bankq}
\begin{bankq}[id=SA05-P1-CHEM-S1-Q08, source={Score Advanced - Test 05 (15 Jul 2021)}, paper=1, qno=8, type=single, subtopic={Enthalpies}, answer={C}, answer_src=key, difficulty=5, diff_parts={K2 S1 I2 C1 F0}, diff_note={concepts: energy level diagram, enthalpies | steps: X and Y are formation and combustion enthalpy of glucose}]
Energy level diagram of formation and combustion of components in standard state is given-\par
\bfig{105.8}{SA05_p1_chem_s1_q08_1.png}
Magnitude of X + Y is equal to -
\begin{qoptions}
\opt{A}{$\Delta$H$^0_{\text{c}}$(C) + $\Delta$H$^0_{\text{c}}$(H$_2$)}
\opt{B}{$\Delta$H$^0_{\text{f}}$(CO$_2$) + $\Delta$H$^0_{\text{f}}$(H$_2$O)}
\opt{C}{$\Delta$H$^0_{\text{f}}$(glucose) + $\Delta$H$^0_{\text{c}}$(glucose)}
\opt{D}{$\Delta$H$^0_{\text{f}}$(glucose) + $\Delta$H$^0_{\text{f}}$(CO$_2$) + $\Delta$H$^0_{\text{f}}$(H$_2$O)}
\end{qoptions}
\end{bankq}
\begin{bankpassage}[id=SA39-P1-CHEM-S3-G5-6, source={Score Advanced - Test 39 (20 Apr 2025)}, questions={5-6}]
0.06 moles of an ideal monoatomic gas undergoes a process as shown in figure. The temperature at point `A' is 100 K.\par
(Use : R = $\dfrac{25}3$J/mol.K, 1L.bar = 100J )\par
\bfig{103.1}{SA39_p1_chem_s3_para5_6_1.png}
\end{bankpassage}
\begin{bankq}[id=SA39-P1-CHEM-S3-Q05, source={Score Advanced - Test 39 (20 Apr 2025)}, paper=1, qno=5, type=numerical, subtopic={Enthalpies}, answer={181.25}, answer_src=key, difficulty=8, diff_parts={K3 S1 I2 C2 F1}, diff_note={concepts: standard problem in this chapter, enthalpies, apply the relevant results | steps: apply the relevant results; compute}, group=SA39-P1-CHEM-S3-G5-6]
The magnitude of work done (in Joule) in cyclic process ABCD is
\end{bankq}
\begin{bankq}[id=SA39-P1-CHEM-S3-Q06, source={Score Advanced - Test 39 (20 Apr 2025)}, paper=1, qno=6, type=numerical, subtopic={Enthalpies}, answer={8.75}, answer_src=key, difficulty=8, diff_parts={K3 S1 I2 C2 F1}, diff_note={concepts: standard problem in this chapter, enthalpies, apply the relevant results | steps: apply the relevant results; compute}, group=SA39-P1-CHEM-S3-G5-6]
The magnitude of enthalpy change (in L.bar) in process CD is:
\end{bankq}
\begin{bankpassage}[id=SA39-P2-CHEM-S3-G3-4, source={Score Advanced - Test 39 (20 Apr 2025)}, questions={3-4}]
Consider the following sequence of reactions.\par
(i) PbO$_2$ + conc. HNO$_3\to$ Pb(NO$_3$)$_2$ + H$_2$O + \underline{X} (g)\par
(ii) 2NaHCO$_3\xrightarrow{\Delta}$ Na$_2$CO$_3$ + H$_2$O + \underline{Y} (g)\par
(iii) FeS$_{2(8)}$ + X(g) $\to$ Fe$_2$O$_3$ + \underline{Z} (g)
\end{bankpassage}
\begin{bankq}[id=SA39-P2-CHEM-S3-Q03, source={Score Advanced - Test 39 (20 Apr 2025)}, paper=2, qno=3, type=numerical, subtopic={Enthalpies}, answer={172.4}, answer_src=solved, working={Solved by hand (key missing): result 172.4.}, difficulty=8, diff_parts={K3 S1 I2 C2 F1}, diff_note={concepts: standard problem in this chapter, enthalpies, apply the relevant results | steps: apply the relevant results; compute}, group=SA39-P2-CHEM-S3-G3-4]
Assume ideal behavior of all three gases, 0.5 moles of each X and Y gas are taken in a container and expanded reversibly and adiabatically from 1 dm$^3$ to 4 dm$^3$ starting from 27$^\circ$C. Find final temperature (K).\par
Given (4)$^{0.4}$ = 1.74
\end{bankq}
\begin{bankq}[id=SA39-P2-CHEM-S3-Q04, source={Score Advanced - Test 39 (20 Apr 2025)}, paper=2, qno=4, type=numerical, subtopic={Enthalpies}, answer={1.15}, answer_src=solved, working={Solved by hand (key missing): result 1.15.}, difficulty=8, diff_parts={K3 S1 I2 C2 F1}, diff_note={concepts: standard problem in this chapter, gas laws, apply the relevant results | steps: apply the relevant results; compute}, group=SA39-P2-CHEM-S3-G3-4]
If equal moles of X,Y,Z (assuming ideal behavior of the gas) are expanded adiabatically and reversibly from the same initial state to same final volume. The magnitude of work is maximum for a gas. What is the value of $\gamma$ for that gas, if all the degree of freedom are active.
\end{bankq}
\begin{bankq}[id=SA40-P1-CHEM-S1-Q03, source={Score Advanced - Test 40 (27 Apr 2025)}, paper=1, qno=3, type=multi, subtopic={Enthalpies}, answer={C,D}, answer_src=key, difficulty=8, diff_parts={K3 S1 I2 C2 F1}, diff_note={concepts: standard problem in this chapter, enthalpies, apply the relevant results | steps: apply the relevant results; compute}]
Among the following solutions, which two solutions will show `highest' and `lowest' increase in temperature on mixing?\par
Assume that initially all solutions at 25$^\circ$C and density of solutions are identical.
\begin{qoptions}
\opt{A}{0.4 M, 50 mL NaOH with 0.5 M, 100 mL HNO$_3$}
\opt{B}{0.2 M, 75 mL H$_2$SO$_4$ with 0.6 M, 25 mL NaOH}
\opt{C}{0.3 M, 20 mL HBr with 0.3 M, 80 mL Ba(OH)$_2$}
\opt{D}{0.9 M, 40 mL Ba(OH)$_2$ with 0.7 M, 60 mL HI}
\end{qoptions}
\end{bankq}

\section*{SRGBT papers}
\begin{bankq}[id=SRGBT04-P1-CHEM-S1-Q20, source={SRGBT - Test 04 (15 Aug 2021)}, paper=1, qno=20, type=multi, subtopic={Enthalpies}, answer={D}, answer_src=key, difficulty=5, diff_parts={K2 S1 I1 C1 F1}, diff_note={concepts: standard problem in this chapter, enthalpies | steps: apply the relevant results; compute}]
Which of the following option is incorrect ?
\begin{qoptions}
\opt{A}{The difference between ionisation enthalpy ($\Delta$H$_\mathrm{IE}$) and ionisation energy (IE) is directly depencent on absolute temperature.}
\opt{B}{heat added to the system at lower temperature causes greater randomness as compared to same amount of heat added at higher temperature.}
\opt{C}{A container having N$_2$O$_4$ gas if heated will show an increase in intensity of brown colour.}
\opt{D}{It is easier to remove H$^+$ from HPO$_4^{2-}$ as compared to that from H$_2$PO$_4$.}
\end{qoptions}
\end{bankq}
\begin{bankq}[id=SRGBT04-P1-CHEM-S1-Q21, source={SRGBT - Test 04 (15 Aug 2021)}, paper=1, qno=21, type=multi, subtopic={Hess's law}, answer={A}, answer_src=key, difficulty=5, diff_parts={K2 S1 I1 C1 F1}, diff_note={concepts: standard problem in this chapter, hess's law | steps: apply the relevant results; compute}]
From the following data.\par
C(s, graphite) + O$_2$(g) $\longrightarrow$ CO$_2$(g) \hfill $\Delta$H$^\circ_\mathrm{reaction}$ = $-$393.5 kJ / mole\par
H$_2$(g) + $\dfrac12$O$_2$(g) $\longrightarrow$ H$_2$O($\ell$) \hfill $\Delta$H$^\circ_\mathrm{reaction}$ = $-$286 kJ / mole\par
2C$_2$H$_6$(g) + 7O$_2$(g) $\longrightarrow$ 4CO$_2$(g) + 6H$_2$O($\ell$) \hfill $\Delta$H$^\circ_\mathrm{reaction}$ = $-$3120 kJ / mole\par
Calculate the standard enthalpy of formation of C$_2$H$_6$ (g) (in kJ/mole).
\begin{qoptions}
\opt{A}{$-$85}
\opt{B}{$-$170}
\opt{C}{170}
\opt{D}{85}
\end{qoptions}
\end{bankq}
\begin{bankq}[id=SRGBT04-P1-CHEM-S1-Q27, source={SRGBT - Test 04 (15 Aug 2021)}, paper=1, qno=27, type=multi, subtopic={Enthalpies}, answer={B,D}, answer_src=key, difficulty=8, diff_parts={K3 S1 I2 C2 F1}, diff_note={concepts: standard problem in this chapter, enthalpies, apply the relevant results | steps: apply the relevant results; compute}]
Select the correct option(s) :
\begin{qoptions}
\opt{A}{$\Delta$H$^\circ_\mathrm f$[H$_2$O($\ell$)] is zero at 25$^\circ$C}
\opt{B}{$\Delta$H$^\circ_\mathrm f$[O$_2$(g)] is zero at 25$^\circ$}
\opt{C}{S$^\circ$[O$_2$(g)] is zero at 25$^\circ$C}
\opt{D}{Entropies of aqueous ions can be negative or positive.}
\end{qoptions}
\end{bankq}
\begin{bankq}[id=SRGBT04-P1-CHEM-S1-Q31, source={SRGBT - Test 04 (15 Aug 2021)}, paper=1, qno=31, type=multi, subtopic={Enthalpies}, answer={A,B,C,D}, answer_src=key, difficulty=8, diff_parts={K3 S1 I2 C2 F1}, diff_note={concepts: standard problem in this chapter, enthalpies, apply the relevant results | steps: apply the relevant results; compute}]
Which is correct relation between various enthalpies given all enthalpies are measured under similar pressure and temperature conditions-
\begin{qoptions}
\opt{A}{2$\Delta_\mathrm f$H$^\circ$(HCl,g) = BE$_\mathrm{(H-H)}$ + BE$_\mathrm{(Cl-Cl)}$ $-$ 2BE$_\mathrm{(H-Cl)}$}
\opt{B}{$\Delta_\mathrm f$H$^\circ$(CO$_2$,g) = $\Delta$H$^\circ_\mathrm{comb.}$ (C, graphite)}
\opt{C}{$\Delta$H$^\circ$ for conversion of d-lactic acid to $l$-lactic acid is zero}
\opt{D}{$\Delta$H$^\circ_\mathrm{neutralization}$ (HCl + NaOH) is equal to $\Delta$H$^\circ$ for the reaction H$^+$(aq.) + OH$^-$(aq.) $\to$ H$_2$O($l$)}
\end{qoptions}
\end{bankq}
\begin{bankq}[id=SRGBT04-P1-CHEM-S1-Q33, source={SRGBT - Test 04 (15 Aug 2021)}, paper=1, qno=33, type=multi, subtopic={Bond energies}, answer={85}, answer_src=key, difficulty=8, diff_parts={K3 S1 I2 C2 F1}, diff_note={concepts: standard problem in this chapter, bond energies, apply the relevant results | steps: apply the relevant results; compute}]
Calculate $\Delta$H$_\mathrm{resonance\;energy}$ of pyridine.\par
Given data : $\Delta$H$^\circ_\mathrm f$ pyridine$_{(\ell)}$ = $-$175 kJ/mole\par
\bfig{325.6}{SRGBT04_p1_chem_s0_q33_1.png}
\end{bankq}

\section*{Enthusiast Pre-Score papers}
\begin{bankq}[id=EPS01-P2-CHEM-S1-Q06, source={Enthusiast Pre-Score - Test 01 (04 Nov 2020)}, paper=2, qno=6, type=multi, subtopic={Enthalpies}, answer={A,B}, answer_src=key, difficulty=7, diff_parts={K2 S2 I2 C1 F1}, diff_note={concepts: hydration and solution enthalpies, Born Haber cycle | steps: Born Haber cycle; enthalpy of hydration; enthalpy of solution}]
From the following data, select the correct option(s)\par
Given : $(\Delta_fH^0)_{\mathrm{AlF_3(s)}}=-1500$ kJ/mol , $(\Delta_fH^0)_{\mathrm{Al^{+3}(aq.)}}=-530$kJ/mol\par
$(\Delta_fH^0)_{\mathrm{F^-(aq.)}}=-330$ kJ/mol , $(\Delta_fH^0)_{\mathrm{Al^{+3}(g)}}=+5480$ kJ/mol
\begin{qoptions}
\opt{A}{$(\Delta H_{\mathrm{hydration}})_{\mathrm{Al^{+3}(g)}}=-6010$ kJ/mol}
\opt{B}{$(\Delta H_{\mathrm{solution}})_{\mathrm{AlF_3(s)}}=-20$ kJ/mol}
\opt{C}{$(\Delta H_{\mathrm{solution}})_{\mathrm{AlF_3(s)}}=+20$ kJ/mol}
\opt{D}{$(\Delta H_{\mathrm{hydration}})_{\mathrm{Al^{+3}(g)}}=-4950$ kJ/mol}
\end{qoptions}
\end{bankq}
\begin{bankq}[id=EPS01-P2-CHEM-S2-Q01, source={Enthusiast Pre-Score - Test 01 (04 Nov 2020)}, paper=2, qno=1, type=single, subtopic={Bond energies}, answer={54.50}, answer_src=key, difficulty=3, diff_parts={K1 S2 I1 C0 F0}, diff_note={concepts: reaction enthalpy from bond energies | steps: bonds broken; bonds formed; sum}]
If $\epsilon_{\mathrm{H-Br}}=368$ kJ/mol, $\epsilon_{\mathrm{H-H}}=435$ kJ/mol , $\epsilon_{\mathrm{Br-Br}}=192$ kJ/mol , then find magnitude of $\Delta_rH^0$(kJ / mol) for the reaction :\par
$\dfrac12$H$_2$(g) $+$ $\dfrac12$Br$_2$(g) $\to$ HBr(g)
\end{bankq}
\begin{bankq}[id=EPS01-P2-CHEM-S2-Q02, source={Enthusiast Pre-Score - Test 01 (04 Nov 2020)}, paper=2, qno=2, type=single, subtopic={Enthalpies}, answer={1318.00}, answer_src=key, difficulty=3, diff_parts={K1 S2 I1 C0 F0}, diff_note={concepts: ionisation enthalpy from a formation enthalpy | steps: atomisation; ionisation; relation}]
If $(\Delta_fH^\circ)_{\mathrm{H^+(g)}}=1536$ kJ/mol and enthalpy of atomisation of H$_2$(g) is 436 kJ/mol. then ionisation enthalpy of H atom in kJ/mol is :
\end{bankq}

\section*{TOAS papers}
\begin{bankq}[id=TOAS10-P1-CHEM-S3-Q03, source={TOAS - Test 10 (12 Jun 2022)}, paper=1, qno=3, type=numerical, subtopic={Hess's law}, answer={961.00}, answer_src=key, difficulty=2, diff_parts={K1 S1 I0 C0 F1}, diff_note={concepts: Hess's law from combustion data | steps: formation = 3 times combustion of C + 4 times that of H2 - combustion of C3H8}]
The standard enthalpies of combustion of C(graphite), H$_2$(g) and C$_3$H$_8$(g) are $-393$, $-285$ and $-3280$ kJ/mole, respectively. The standard enthalpy of formation (in KJ/mole) of C$_3$H$_8$(g) is.
\end{bankq}
\begin{bankq}[id=TOAS11-P2-CHEM-S1-Q02, source={TOAS - Test 11 (27 Jun 2022)}, paper=2, qno=2, type=multi, subtopic={Enthalpies}, answer={A,B,C}, answer_src=key, difficulty=2, diff_parts={K1 S1 I0 C0 F1}, diff_note={concepts: standard enthalpy of formation | steps: reaction must form one mole of a compound from elements in standard states; check each}]
Which of the following reaction(s) represent standard enthalpy of formation of the product formed ?
\begin{qoptions}
\opt{A}{C(graphite) $\to$ C(diamond)}
\opt{B}{Hg ($\ell$) $\to$ Hg(g)}
\opt{C}{$\frac12$ I$_2$(solid) $+$ $\frac12$Cl$_2$(g) $\to$ ICl(s)}
\opt{D}{$\frac12$ Br$_2$(g) $+$ $\frac12$O$_2$(g) $\to$ BrO(g)}
\end{qoptions}
\end{bankq}
\begin{bankpassage}[id=TOAS12-P2-CHEM-S2-G11-12, source={TOAS - Test 12 (10 Jul 2022)}, questions={11-12}]
Enthalpy of combustion of a given compounds is the change in enthalpy when 1 mole of this compound combines with required amount of O$_2$ to give product in their stable forms. The enthalpy of combustion is usually measured by placing a known mass of the compound in a closed steel container (known as bomb-calorimeter). Graphite is thermodynamically most stable form of carbon.
\end{bankpassage}
\begin{bankq}[id=TOAS12-P2-CHEM-S2-Q11, source={TOAS - Test 12 (10 Jul 2022)}, paper=2, qno=11, type=multi, subtopic={Enthalpies}, answer={A,B,C}, answer_src=key, difficulty=4, diff_parts={K1 S2 I0 C1 F1}, diff_note={concepts: bond enthalpy from combustion data | steps: Hess cycle for methane atomisation; C-H bond enthalpy; check each statement}, group=TOAS12-P2-CHEM-S2-G11-12]
Given the data below :\par
$\Delta H^0_{\text{comb}}$(CH$_4$, g) $=-900$ kJ/mole\par
$\Delta H^0_{\text{comb}}$(C, graphite) $=-400$ kJ/mole\par
$\Delta H^0_{\text{comb}}$(H$_2$, g) $=-300$ kJ/mole\par
$\Delta H^0_{\text{atomization}}$(C, graphite) $=700$ kJ/mole\par
and $\Delta H^0_f$(H, g) $=220$ kJ/mole\par
Identify the incorrect option(s) :
\begin{qoptions}
\opt{A}{Bond enthalpy of C--H bond is 340 kJ/mole}
\opt{B}{1 mole graphite produces more heat on combustion than 1 mole diamond at 1 atm and 25$^\circ$C}
\opt{C}{Enthalpy of atomisation of CH$_4$(g) is 1600 kJ/mole}
\opt{D}{Bond enthalpy of C--H bond is 420 kJ/mole}
\end{qoptions}
\end{bankq}
\begin{bankq}[id=TOAS12-P2-CHEM-S2-Q12, source={TOAS - Test 12 (10 Jul 2022)}, paper=2, qno=12, type=multi, subtopic={Enthalpies}, answer={A,B,C}, answer_src=key, difficulty=3, diff_parts={K1 S2 I0 C0 F1}, diff_note={concepts: definitions of enthalpy of formation and combustion | steps: C graphite combustion equals CO2 formation; H2 combustion equals H2O formation; signs}, group=TOAS12-P2-CHEM-S2-G11-12]
Identify the correct statement -
\begin{qoptions}
\opt{A}{$\Delta H^0_{\text{comb}}$(C, graphite) is equal to $\Delta H^0_f$(CO$_2$, g)}
\opt{B}{$\Delta H^0_f$(H$_2$O, $\ell$) is equal to $\Delta H^0_{\text{comb}}$(H$_2$, g)}
\opt{C}{The standard enthalpy of combustion is always negative.}
\opt{D}{The standard enthalpy of formation is always negative}
\end{qoptions}
\end{bankq}
\begin{bankq}[id=TOAS13-P2-CHEM-S1-Q01, source={TOAS - Test 13 (24 Jul 2022)}, paper=2, qno=1, type=single, subtopic={Hess's law}, answer={D}, answer_src=key, difficulty=2, diff_parts={K1 S2 I0 C0 F0}, diff_note={concepts: Hess law | steps: reverse the first equation and halve it; double the second; add}]
Joseph Priestley discovered sweet smelling, colorless laughing gas, nitrous oxide (N$_2$O), in 1763. N$_2$O has been widely used as an anesthetic for medical and dental applications. Nitrogen dioxide (NO$_2$) is a toxic reddish brown gas that is the major component of smog. Nitrogen oxide (NO) is a colorless, non-toxic gas that readily oxidizes to nitrogen dioxide in the presence of oxygen. Using the reactions :\par
4NO(g) $\to$ 2N$_2$O(g) $+$ O$_2$(g) \quad $\Delta H^0=-139.56$ kJ\par
2NO(g) $+$ O$_2$(g) $\to$ 2NO$_2$(g) \quad $\Delta H^0=-69.70$ kJ\par
Calculate the standard enthalpy change when 2 moles of nitrous oxide combine with three moles of oxygen to produce four moles of nitrogen dioxide.
\begin{qoptions}
\opt{A}{$-209.26$ kJ}
\opt{B}{$-139.86$ kJ}
\opt{C}{$-69.86$ kJ}
\opt{D}{$+0.16$ kJ}
\end{qoptions}
\end{bankq}
\begin{bankq}[id=TOAS13-P2-CHEM-S1-Q05, source={TOAS - Test 13 (24 Jul 2022)}, paper=2, qno=5, type=single, subtopic={Enthalpies}, answer={B}, answer_src=key, difficulty=1, diff_parts={K1 S1 I0 C0 F0}, diff_note={concepts: heat of neutralisation | steps: heat depends on moles of the limiting reagent and the volume; weak acids absorb ionisation heat so give a smaller rise}]
When 5 mL of a 1.0 M HCl solution is mixed with 5 mL of a 0.1 M NaOH solution, temperature of solution increases by $t^\circ$C. Which of the following can be predicted from this observation :
\begin{qoptions}
\opt{A}{If 5 mL of 0.1M HCl is mixed with 5 mL 0.1 M NH$_3$ solution, the temperature rise will be more than $t^\circ$C}
\opt{B}{If 5 mL of 0.1 M CH$_3$COOH is mixed with 5 mL 0.1 M NaOH, the temperature rise will be less than $t^\circ$C}
\opt{C}{If 10 mL of 1.0 M HCl mixed with 10 mL of 0.1 M NaOH, the temperature rise will be $2t^\circ$C}
\opt{D}{If 10 mL of 0.1 M HCl is mixed with 10 mL of 0.1 M NaCl, the temperature rise will be $t^\circ$C}
\end{qoptions}
\end{bankq}
\begin{bankpassage}[id=TOAS14-P1-CHEM-S2-G1-2, source={TOAS - Test 14 (07 Aug 2022)}, questions={1-2}]
When 100 mL of 1.0 M HCl was mixed with 100 mL of 1.0 M NaOH in an insulated beaker at constant pressure, a temperature increase of 5.7 $^\circ$C was measured for the beaker and its contents.\par
(Expt-1). Because the enthalpy of neutralisation of a strong acid with a strong base is a constant ($-57.0$ kJmol$^{-1}$), this experiment could be used to measure the calorimeter constant. In a second experiment (Expt-2), 100 mL of 2.0 M acetic acid ($K_a=2.0\times10^{-5}$) was mixed with 100 mL of 1.5 M NaOH (under identical conditions to (Expt-1)) where a temperature rise of 8.25 $^\circ$C was measured.\par
(Consider heat capacity of all solutions as 4.2 Jg$^{-1}$K$^{-1}$ and density of all solutions as 1.0 g mL$^{-1}$)
\end{bankpassage}
\begin{bankq}[id=TOAS14-P1-CHEM-S2-Q01, source={TOAS - Test 14 (07 Aug 2022)}, paper=1, qno=1, type=numerical, subtopic={Enthalpies}, answer={2.00}, answer_src=key, difficulty=5, diff_parts={K2 S2 I1 C0 F1}, diff_note={concepts: calorimeter constant, enthalpy of dissociation of a weak acid | steps: heat of expt 1 gives the total heat capacity; heat in expt 2 divided by moles of base gives the neutralisation enthalpy; subtract the strong acid value}, group=TOAS14-P1-CHEM-S2-G1-2]
Enthalpy of dissociation (in kJ mol$^{-1}$) of acetic acid obtained from the Expt-2 is
\end{bankq}
\begin{bankq}[id=TOAS14-P1-CHEM-S2-Q02, source={TOAS - Test 14 (07 Aug 2022)}, paper=1, qno=2, type=numerical, subtopic={Enthalpies}, answer={5.18}, answer_src=key, difficulty=4, diff_parts={K1 S2 I0 C1 F1}, diff_note={concepts: buffer pH from Henderson equation | steps: moles of acid and acetate after reaction; pKa from Ka; log ratio}, group=TOAS14-P1-CHEM-S2-G1-2]
The pH of the solution after Expt-2
\end{bankq}
\begin{bankpassage}[id=TOAS16-P2-CHEM-S3-G1-2, source={TOAS - Test 16 (04 Sep 2022)}, questions={1-2}]
Consider the combustion of 1 mole of CH$_4$ according to the following reaction with 20\% excess air at $27^\circ$C and 1 bar pressure.\par
\noindent\hfill CH$_4$(g) + 2O$_2$(g) $\to$ CO$_2$(g) + 2H$_2$O(g)\hfill\null\par
Assume air as a mixture of 20\% O$_2$ and 80\% N$_2$ by volume.\par
Given : $\Delta H^\circ_f$, CO$_2$(g) $=-100$ K cal mol$^{-1}$\par
$\Delta H^\circ_f$, H$_2$O(g) $=-60$ K cal mol$^{-1}$\par
$\Delta H^\circ_f$, CH$_4$(g) $=-20$ K cal mol$^{-1}$\par
Assume all the gases are ideal gases and vibrational degree of freedom is not active. $R=2$ cal K$^{-1}$ mol$^{-1}$.
\end{bankpassage}
\begin{bankq}[id=TOAS16-P2-CHEM-S3-Q01, source={TOAS - Test 16 (04 Sep 2022)}, paper=2, qno=1, type=numerical, subtopic={Enthalpies}, answer={2450.53 or 2450.54}, answer_src=key, difficulty=7, diff_parts={K3 S2 I1 C1 F1}, diff_note={concepts: enthalpy of combustion, heat capacities, adiabatic flame temperature | steps: delta H of the reaction; heat capacities of each product; sum for the products and the excess air; temperature rise}, group=TOAS16-P2-CHEM-S3-G1-2]
Calculate maximum temperature in Kelvin that can be reached by above combustion. Ignore the variation of $\Delta_rH$ with temperature.
\end{bankq}
\begin{bankq}[id=TOAS16-P2-CHEM-S3-Q02, source={TOAS - Test 16 (04 Sep 2022)}, paper=2, qno=2, type=numerical, subtopic={Enthalpies}, answer={159.92}, answer_src=key, difficulty=7, diff_parts={K3 S2 I1 C1 F1}, diff_note={concepts: enthalpy at a different temperature, Carnot efficiency | steps: delta H at 127 C by Kirchhoff; maximum efficiency; 80 percent of it; heat discarded}, group=TOAS16-P2-CHEM-S3-G1-2]
If energy released by combustion of 1 mole of CH$_4$ at $127^\circ$C and 1 bar is supplied to a heat engine at $127^\circ$C which discards heat at $27^\circ$C and whose efficiency is 80\% of maximum possible efficiency then calculate amount of heat discarded by the heat engine in Kcal. ($R=2$ cal K$^{-1}$ mol$^{-1}$).
\end{bankq}
\begin{bankq}[id=TOAS17-P1-CHEM-S3-Q02, source={TOAS - Test 17 (18 Sep 2022)}, paper=1, qno=2, type=single, subtopic={Bond energies}, answer={0.78}, answer_src=key, difficulty=6, diff_parts={K2 S3 I1 C1 F0}, diff_note={concepts: enthalpy of formation and atomisation, bond energies | steps: atomisation of carbon and hydrogen; formation cycle; C-H contribution; bond energy; scale}]
Using the data provided, calculate the multiple bond energy (kJ mol$^{-1}$) of a C$\equiv$C in C$_2$H$_2$. That energy is (take the bond energy of C--H bond as 350 kJ bond mol$^{-1}$).\par
2C(s) + H$_2$(g) $\to$ C$_2$H$_2$(g) \ ; $\Delta H=220$ kJ/mol\par
2C(s) $\to$ 2C(g) \ ; $\Delta H=1400$ kJ/mol\par
H$_2$(g) $\to$ 2H(g) \ ; $\Delta H=300$ kJ/mol\par
If the value of bond enthalpy is $x$ kJ/mol ; calculate $\left(\dfrac x{1000}\right)$
\end{bankq}
\begin{bankq}[id=TOAS17-P2-CHEM-S1-Q03, source={TOAS - Test 17 (18 Sep 2022)}, paper=2, qno=3, type=single, subtopic={Enthalpies}, answer={B}, answer_src=key, difficulty=1, diff_parts={K1 S1 I0 C0 F0}, diff_note={concepts: enthalpy of neutralisation of a weak acid | steps: ionisation enthalpy = observed neutralisation - strong acid neutralisation per mole of H+; two H+ ions}]
Enthalpy of neutralization of H$_3$PO$_3$ acid is $-106.68$ kJ/mol using NaOH. Enthalpy of neutralization of HCl by NaOH is $-55.84$ kJ/mol. Calculate $\Delta H_{\text{ionization}}$ of H$_3$PO$_3$ into its ions:
\begin{qoptions}
\opt{A}{50.84 kJ/mol}
\opt{B}{5 kJ/mol}
\opt{C}{2.5 kJ/mol}
\opt{D}{None of these}
\end{qoptions}
\end{bankq}
\begin{bankq}[id=TOAS17-P2-CHEM-S1-Q04, source={TOAS - Test 17 (18 Sep 2022)}, paper=2, qno=4, type=single, subtopic={Bond energies}, answer={A}, answer_src=key, difficulty=4, diff_parts={K2 S2 I0 C1 F0}, diff_note={concepts: bond energies and resonance energy, vaporisation | steps: bonds broken and formed in propane combustion; vaporisation of water; resonance correction}]
The enthalpy of combustion of propane (C$_3$H$_8$) gas in terms of given data is :
\par\smallskip\noindent Bond energy (kJ/mol)\par\smallskip
\begin{center}\footnotesize\setlength{\tabcolsep}{2pt}
\begin{tabular}{|c|c|c|c|c|}\hline
$\varepsilon_{\text{C--H}}$&$\varepsilon_{\text{O=O}}$&$\varepsilon_{\text{C=O}}$&$\varepsilon_{\text{O--H}}$&$\varepsilon_{\text{C--C}}$\\\hline
$+x_1$&$+x_2$&$+x_3$&$+x_4$&$+x_5$\\\hline
\end{tabular}\end{center}
Resonance energy of CO$_2$ is $-\,z$ kJ/mol and $\Delta H_{\text{vaporization}}$ [H$_2$O($\ell$)] is $y$ kJ/mol.
\begin{qoptions}
\opt{A}{$8x_1+2x_5+5x_2-6x_3-8x_4-4y-3z$}
\opt{B}{$6x_1+x_5+5x_2-3x_3-4x_4-4y-3z$}
\opt{C}{$8x_1+2x_5+5x_2-6x_3-8x_4-y-z$}
\opt{D}{$8x_1+x_5+5x_2-6x_3-8x_4-4y+3z$}
\end{qoptions}
\end{bankq}
\begin{bankq}[id=TOAS18-P2-CHEM-S1-Q02, source={TOAS - Test 18 (30 Oct 2022)}, paper=2, qno=2, type=single, subtopic={Enthalpies}, answer={A}, answer_src=key, difficulty=2, diff_parts={K1 S2 I0 C0 F0}, diff_note={concepts: calorimetry of neutralisation | steps: equivalents reacting; heat; temperature rise}]
In a constant pressure calorimeter 224 ml of 0.1M KOH (aq.) solution is reacted with 50 ml of 0.1M H$_2$SO$_4$ (aq.) solution then increase in temperature of solution will be (assume heat capacity of calorimeter is negligible) -
\par\noindent Given :\par
\noindent Specific heat of solution = 1 cal/gm--K\par
\noindent Density of solution = 1 gm/ml\par
\noindent HCl(aq.) + NaOH(aq.) $\to$ NaCl (aq.) + H$_2$O(l) $\Delta H=-13.7$ kcal/eq.
\begin{qoptions}
\opt{A}{0.5 K}
\opt{B}{1 K}
\opt{C}{2K}
\opt{D}{2.4 K}
\end{qoptions}
\end{bankq}

\section*{TAAS papers}
\begin{bankq}[id=TAAS02-P1-CHEM-S2-Q07, source={TAAS - Test 02 (30 Jul 2023)}, paper=1, qno=7, type=multi, subtopic={Enthalpies}, answer={A,B,C}, answer_src=key, difficulty=3, diff_parts={K1 S2 I0 C0 F1}, diff_note={concepts: enthalpy of formation and combustion and atomisation | steps: definition of each enthalpy; match the reaction; check each statement}, also_in={TNAS - Test 06 (30 Jul 2023) P1 Q6}]
Which is/are correct statement -
\begin{qoptions}
\opt{A}{$\mathrm{H_2(g)+\tfrac12O_2(g)\to H_2O(l)}$ ; $\Delta_r H$

$\Delta_rH=\Delta H_C(\mathrm{H_2,g})=\Delta H_f(\mathrm{H_2O},\ell)$}
\opt{B}{Na(s)$\to$ Na(g) ; $\Delta_rH$

$\Delta_rH=\Delta H_{\text{sublimation}}(\mathrm{Na,s})=\Delta H_{\text{atomisation}}(\mathrm{Na,s})$}
\opt{C}{C(graphite) $+$ O$_2$(g)$\to$ CO$_2$(g) ; $\Delta_rH$

$\Delta_rH=\Delta H_C(\mathrm{C,graphite})=\Delta H_f(\mathrm{CO_2,g})$}
\opt{D}{H$_2$(g)$\to$ 2H(g) ; $\Delta_rH$

$\Delta_rH=\Delta H_{\text{atomisation}}(\mathrm{H_2,g})=\Delta H_f(\mathrm{H,g})$}
\end{qoptions}
\end{bankq}

\section*{TOPS papers}
\begin{bankq}[id=TOPS01-P1-CHEM-S1-Q03, source={TOPS - Test 01 (31 Jul 2022)}, paper=1, qno=3, type=multi, subtopic={Enthalpies}, answer={B,C}, answer_src=key, difficulty=5, diff_parts={K2 S2 I0 C1 F1}, diff_note={concepts: enthalpy of formation, lattice and hydration energy | steps: standard states; sign of hydration; definition of lattice energy}]
Consider the following statements :

I. Enthalpy of formation is always calculated from its constituent elements in their solid state.

II. Hydration energy of ion is always amount of energy released.

III. Lattice energy is amount of energy supplied to an ionic solid so that it is melted into constituent ions
\begin{qoptions}
\opt{A}{I \& II are correct}
\opt{B}{II is correct}
\opt{C}{I \& III are incorrect}
\opt{D}{I, II \& III are incorrect}
\end{qoptions}
\end{bankq}

\section*{TEAS papers}
\begin{bankpassage}[id=TEAS01-P1-CHEM-S2-G9-10, source={TEAS - Test 01 (16 Apr 2023)}, questions={9-10}]
Study the following thermochemical equations :\\
(a) $\mathrm{N_2O(g)+3H_2(g)\to H_2O(\ell)+N_2H_4(\ell)}$; $\Delta H=-76$ Kcal\\
(b) $\mathrm{4N_2(g)+3H_2O(\ell)\to 2NH_3(g)+3N_2O(g)}$; $\Delta H=+240$ Kcal\\
(c) $\mathrm{N_2H_4(\ell)+H_2O(\ell)\to 2NH_3(g)+\tfrac12O_2(g)}$; $\Delta H=+36$ Kcal\\
(d) $\mathrm{H_2O(\ell)\to H_2(g)+\tfrac12O_2(g)}$; $\Delta H=+68$ Kcal
\end{bankpassage}
\begin{bankq}[id=TEAS01-P1-CHEM-S2-Q09, source={TEAS - Test 01 (16 Apr 2023)}, paper=1, qno=9, type=single, subtopic={Hess's law}, answer={C}, answer_src=key, difficulty=6, diff_parts={K2 S2 I2 C1 F0}, diff_note={concepts: Hess law, enthalpy of combustion | steps: combine the equations; enthalpy of formation of water or hydrazine; combustion reaction; value}, group=TEAS01-P1-CHEM-S2-G9-10]
The enthalpy of combustion (in Kcal/mol) of liquid $\mathrm{N_2H_4}$ is :
\begin{qoptions}
\opt{A}{$-129$}
\opt{B}{$-138$}
\opt{C}{$-147$}
\opt{D}{$+147$}
\end{qoptions}
\end{bankq}
\begin{bankq}[id=TEAS01-P1-CHEM-S2-Q10, source={TEAS - Test 01 (16 Apr 2023)}, paper=1, qno=10, type=single, subtopic={Hess's law}, answer={C}, answer_src=key, difficulty=5, diff_parts={K1 S2 I2 C1 F0}, diff_note={concepts: Hess law | steps: combine the equations; formation reaction; value}, group=TEAS01-P1-CHEM-S2-G9-10]
The enthalpy of formation (in Kcal/mol) of liquid $\mathrm{N_2H_4}$ is :
\begin{qoptions}
\opt{A}{$-11$}
\opt{B}{$+44$}
\opt{C}{$+11$}
\opt{D}{$+13$}
\end{qoptions}
\end{bankq}

\section*{TMAS papers}
\begin{bankq}[id=TMAS02-P1-CHEM-S1-Q02, source={TMAS - Test 02 (06 Aug 2023)}, paper=1, qno=2, type=single, subtopic={Enthalpies}, answer={A}, answer_src=key, difficulty=4, diff_parts={K1 S2 I1 C1 F0}, diff_note={concepts: enthalpy relation | steps: enthalpy of formation relation; ratio condition; solve}]
Enthalpy for the reaction :\\
$\mathrm{Ag^+(aq)+Br^-(aq)\to AgBr(s)}$ is $-100$ KJ.\\
If magnitude of enthalpy of formation of $\mathrm{Ag^+(aq)}$ and $\mathrm{Br^-(aq)}$ is in the ratio of $5:6$ while formation of $\mathrm{Ag^+(aq)}$ is an endothermic where as formation of $\mathrm{Br^-}$ is an exothermic process. If enthalpy of formation of $\mathrm{AgBr(s)}=-110\ \mathrm{KJ\,mol^{-1}}$ then the enthalpy of formation of $\mathrm{Ag^+(aq)}$ is :
\begin{qoptions}
\opt{A}{$50\ \mathrm{KJ\,mol^{-1}}$}
\opt{B}{$60\ \mathrm{KJ\,mol^{-1}}$}
\opt{C}{$1000\ \mathrm{KJ\,mol^{-1}}$}
\opt{D}{$10\ \mathrm{KJ\,mol^{-1}}$}
\end{qoptions}
\end{bankq}

\section*{TNAS papers}
\begin{bankq}[id=TNAS04-P2-CHEM-S3-Q02, source={TNAS - Test 04 (18 Jun 2023)}, paper=2, qno=2, type=numerical, subtopic={Enthalpies}, answer={0.60}, answer_src=key, difficulty=7, diff_parts={K1 S3 I1 C2 F1}, diff_note={concepts: heat balance of two reactions with different rates | steps: rate of the first reaction; heat released; heat absorbed by the mixture; rate of the second reaction; rate of Q}]
For hypothetical reaction occurring in same container\par
$\mathrm{A+8B\longrightarrow P}$\quad $\Delta H=-2x$ KJ\par
$\mathrm{M\longrightarrow 2Q+R}$\quad $\Delta H=x$ KJ\par
Temperature of solution is increased by $2^\circ$C per sec.\par
If rate of disappearance of B = 2 mole sec$^{-1}$\par
Heat capacity of overall reaction mixture = 100 x J/K then what is rate of appearance of Q (in mol/sec).
\end{bankq}
\begin{bankq}[id=TNAS05-P1-CHEM-S1-Q02, source={TNAS - Test 05 (09 Jul 2023)}, paper=1, qno=2, type=single, subtopic={Hess's law}, answer={C}, answer_src=key, difficulty=4, diff_parts={K1 S2 I1 C1 F0}, diff_note={concepts: enthalpy of formation from combustion data | steps: write the formation reaction; combine given equations; evaluate}]
The enthalpies of combustion of S(s) and $\mathrm{H_2(g)}$ are $-300$ and $-290$ kcal/mol respectively.\par
Given :\par
$\mathrm{SO_3(g)+H_2O(\ell)\longrightarrow H_2SO_4(\ell)}$ $\Delta H=-130$ kcal/mol\par
$\mathrm{SO_2(g)+\frac12O_2(g)\longrightarrow SO_3(g)}$ $\Delta H=-100$ kcal/mol\par
The enthalpy of formation of $\mathrm{H_2SO_4(\ell)}$ would be
\begin{qoptions}
\opt{A}{$-230$ kcal/mol}
\opt{B}{$-620$ kcal/mol}
\opt{C}{$-820$ kcal/mol}
\opt{D}{$-560$ kcal/mol}
\end{qoptions}
\end{bankq}
\begin{bankq}[id=TNAS06-P2-CHEM-S1-Q03, source={TNAS - Test 06 (30 Jul 2023)}, paper=2, qno=3, type=single, subtopic={Bond energies}, answer={A}, answer_src=key, difficulty=4, diff_parts={K1 S2 I1 C1 F0}, diff_note={concepts: enthalpy of atomisation and bond enthalpies | steps: structure of the oxide; count of bonds; bond energy total; sublimation}]
The enthalpy of atomization of $\mathrm{P_4O_6(s)}$ is 3500 kJ/mol. If P--O bond enthalpy is 250 kJ/mol then enthalpy of sublimation of $\mathrm{P_4O_6(s)}$ is :
\begin{qoptions}
\opt{A}{500 kJ/mol}
\opt{B}{600 kJ/mol}
\opt{C}{300 kJ/mol}
\opt{D}{200 kJ/mol}
\end{qoptions}
\end{bankq}
\begin{bankq}[id=TNAS07-P1-CHEM-S2-Q05, source={TNAS - Test 07 (20 Aug 2023)}, paper=1, qno=5, type=single, subtopic={Enthalpies}, answer={675}, answer_src=key, difficulty=6, diff_parts={K1 S3 I1 C2 F0}, diff_note={concepts: heat of combustion at constant volume | steps: combustion of butane from formation data; enthalpy; change in moles of gas; internal energy change; magnitude}]
$\Delta_CH^\circ$ of graphite $=-95$ kcal/mol\par
$\Delta_CH^\circ$ of $\mathrm{H_2(g)}=-68$ kcal/mol\par
$\Delta_fH^\circ$ of $\mathrm{C_4H_{10}(g)}=-42.9$ kcal/mol\par
Calculate magnitude of standard heat of combustion of butane gas at constant temperature of $27^\circ$C and constant volume. ($R=2$ cal/mol-K)
\end{bankq}

\section*{Star Review (Other) papers}
\begin{bankpassage}[id=STAR03-P1-CHEM-S3-G1-2, source={Star Review (Other) - Test 03 (04 Dec 2023)}, questions={1-2}]
Amount of heat evolved during complete combustion of liquid benzene can be calculated from the following data.\par
(i) $18$ gm of graphite on complete combustion evolves $591$ kJ heat in standard conditions\par
(ii) $\Delta H_f^\circ(\mathrm{H_2O,\ \ell})=-286$ kJ/mol\par
(iii) The heat of formation of liquid benzene is $50$ kJ /mole in standard conditions.
\end{bankpassage}
\begin{bankq}[id=STAR03-P1-CHEM-S3-Q01, source={Star Review (Other) - Test 03 (04 Dec 2023)}, paper=1, qno=1, type=numerical, subtopic={Enthalpies}, answer={-394.00}, answer_src=key, difficulty=6, diff_parts={K1 S2 I1 C2 F1}, diff_note={concepts: heat of formation of carbon dioxide from graphite data | steps: moles of graphite; heat per mole; formation enthalpy}, group=STAR03-P1-CHEM-S3-G1-2]
Heat of formation of $\mathrm{CO_2(g)}$ from above data (in kJ/mole) in standard conditions is -
\end{bankq}
\begin{bankq}[id=STAR03-P1-CHEM-S3-Q02, source={Star Review (Other) - Test 03 (04 Dec 2023)}, paper=1, qno=2, type=numerical, subtopic={Enthalpies}, answer={3272.00}, answer_src=key, difficulty=7, diff_parts={K1 S2 I2 C2 F1}, diff_note={concepts: heat of combustion of benzene from formation data | steps: enthalpies of formation of the products; combustion enthalpy of benzene; moles in the sample; heat}, group=STAR03-P1-CHEM-S3-G1-2]
Magnitude of heat evolved (in kJ) from combustion of $78$ gm benzene ($\ell$) in standard conditions is :
\end{bankq}
\end{document}
